\section{Model and Procurement Continuation}
\label{sec:model}

\subsection{Regional policy choices}

There are two symmetric regions, indexed by $i\in\{1,2\}$. Each regional government chooses one of two industrial-policy roles,
\begin{equation}
s_i\in\{U,D\}.
\end{equation}
Activity $U$ is the locally attractive lead/buyer-side activity and yields direct real surplus $\Delta>0$ to the region choosing it. Activity $D$ is the complementary supplier-side activity: the region choosing it hosts or develops the supplier base used by the lead side. The direct local return to $D$ apart from procurement-side surplus is normalized to zero. The normalization only fixes the opportunity cost of taking the complementary role: a government that switches from $U$ to $D$ gives up $\Delta$.

The policy roles are technologically complementary across regions. A procurement continuation is activated only when the policy profile is differentiated, $(U,D)$ or $(D,U)$. If both governments choose $U$, each obtains its direct return $\Delta$ and no complementary project occurs. If both choose $D$, neither obtains the direct $U$ return and no complementary project occurs.

A parameter $A>0$ scales the mass or economic importance of the complementary relationship. It can be interpreted as the scale, mass, or economic importance of otherwise identical complementary procurement opportunities. It is a continuous intensity parameter, not literally a count of projects. It does not change the supplier cost distribution.

\subsection{Supplier market}

Conditional on a differentiated policy profile, the $U$ region is the lead/buyer side of the project and procures from a supplier base hosted or developed by the $D$ region. The $D$ region has $N\ge2$ potential suppliers. Supplier $k$ has private production cost
\begin{equation}
C_k\overset{iid}{\sim}F
\end{equation}
on compact support $[c_0,\bar c]$, with continuous strictly positive density $f$. Costs are independent private information and suppliers are risk neutral. The project creates gross value $v$ and the baseline assumes
\begin{equation}
v>\bar c,
\label{eq:guaranteed-trade}
\end{equation}
so trade always occurs once a differentiated policy profile is selected.

For transparency, the canonical implementation is a second-price reverse auction. The lowest-cost supplier wins and is paid the second-lowest cost. Let $C_{1:N}$ and $C_{2:N}$ denote the first and second order statistics. Truthful cost revelation is a dominant strategy, and the realized decomposition is
\begin{align}
\text{real match surplus} &= v-C_{1:N},\\
\text{buyer surplus} &= v-C_{2:N},\\
\text{winning supplier rent} &= C_{2:N}-C_{1:N}.
\end{align}
Define the expected objects
\begin{align}
m_N &\equiv \E[C_{1:N}],\label{eq:mN}\\
s_N &\equiv \E[C_{2:N}],\label{eq:sN}\\
R_N &\equiv s_N-m_N=\E[C_{2:N}-C_{1:N}],\label{eq:RN}\\
B_N &\equiv v-s_N,\label{eq:BN}\\
G_N &\equiv v-m_N=B_N+R_N.\label{eq:GN}
\end{align}
Here $G_N$ is expected real surplus of the active complementary relation, $R_N$ is expected winning supplier information rent, and $B_N$ is expected buyer surplus.

The baseline distribution satisfies decreasing reversed hazard rate (DRHR):
\begin{equation}
\frac{f(c)}{F(c)}\quad\text{is decreasing in }c.
\label{eq:drhr}
\end{equation}
This is a standard regularity condition in procurement models. We use it only for the comparative static of supplier rent. In particular, procurement theory establishes that, under DRHR, expected winning rent $R_N$ is strictly decreasing in $N$ and converges to zero \citep{Li2005,XuLi2019}. That result is imported rather than derived as a contribution of this paper.

\subsection{Regional incidence}

The $U$ region internalizes the buyer-side surplus $B_N$. The $D$ regional government internalizes a fraction
\begin{equation}
\rho\in(0,1]
\end{equation}
of winning supplier rent. Thus its expected local benefit from the procurement continuation is $\rho R_N$. The parameter $\rho$ is an incidence mapping: it is the share or weight of supplier surplus internalized by the $D$ regional government. The remaining $(1-\rho)R_N$ is outside that regional government's objective but, in the baseline, remains within the constituency of the coordinated benchmark. It is therefore not destroyed and is not interpreted mechanically as foreign ownership.

The coordinated benchmark introduced below values the full real surplus $G_N$. Thus $\rho<1$ represents incomplete regional-government internalization within the modeled supralocal constituency, not a social waste term.

\subsection{Timing and induced payoff matrix}

The timing is:
\begin{enumerate}[label=(\arabic*)]
\item $N,F,\Delta,A,v,$ and $\rho$ are common knowledge.
\item Regional governments simultaneously choose $U$ or $D$.
\item If the profile is differentiated, supplier costs are realized privately and the procurement auction is played.
\item Production and payoffs are realized.
\end{enumerate}

Solving the procurement continuation yields the expected first-stage payoff matrix in Table \ref{tab:payoffs}.

\begin{table}[ht]
\centering
\caption{Expected regional payoffs}
\label{tab:payoffs}
\begin{tabular}{c c c}
\toprule
& Region 2: $U$ & Region 2: $D$\\
\midrule
Region 1: $U$ & $(\Delta,\Delta)$ & $(\Delta+A B_N,\; A\rho R_N)$\\
Region 1: $D$ & $(A\rho R_N,\; \Delta+A B_N)$ & $(0,0)$\\
\bottomrule
\end{tabular}
\end{table}

Because $v>\bar c\ge C_{2:N}$ almost surely, $B_N>0$. Hence $U$ is always a strict best response to an opponent choosing $D$. The strategic issue is whether a government is willing to choose $D$ when the other region chooses $U$.
